% Part: applied-modal-logics
% Chapter: epistemic-logic
% Section: bisimulations

\documentclass[../../../include/open-logic-section]{subfiles}

\begin{document}

\olfileid[zh]{aml}{el}{bsd}

\olsection{双模拟}

关于我们认知语言表达能力，还有一个悬而未决的问题涉及模型与在其中为真的公式之间的关系。我们从框架对应结果中已经看到，当某些公式在某个框架中有效时，它们也保证了这些框架满足某些性质。但是，我们的模态语言，例如，是否允许我们区分这样一个世界——其上有一个自反箭头——与一个世界的无穷链，其中每个世界都指向下一个世界呢？也就是说，是否存在某个公式 $A$，它可能只在这两个世界中的一个处为真？

双模拟是我们可以在关系模型之间定义的一种关系，用来表明它们实际上具有相同的结构。而且，正如我们将看到的，它刻画了模型之间某种可用我们认知语言刻画出的等价性。

\begin{defn}[双模拟]
设 $M_1 = \tuple{W_1, R_1, V_1}$ 与 $M_2 = \tuple{W_2, R_2, V_2}$ 是两个关系模型。并设 $\mathcal{R} \subseteq W_1 \times W_2$ 是一个二元关系。我们说，当对每个 $\tuple{w_1, w_2} \in \mathcal{R}$，都有：

\begin{enumerate}
  \item 对所有!!{propositional variable}~$p$，$w_1 \in V_1(p)$ 当且仅当 $w_2 \in V_2(p)$。
  \item 对所有主体 $a \in A$ 与世界 $v_1 \in W_1$，如果 $R_{1_a} w_1 v_1$，则存在某个 $v_2 \in W_2$ 使得 $R_{2_a} w_2 v_2$，且 $\tuple{v_1, v_2} \in \mathcal{R}$。
  \item 对所有主体 $a \in A$ 与世界 $v_2 \in W_2$，如果 $R_{2_a} w_2 v_2$，则存在某个 $v_1 \in W_1$ 使得 $R_{1_a} w_1 v_1$，且 $\tuple{v_1, v_2} \in \mathcal{R}$。
\end{enumerate}

时，$\mathcal{R}$ 就是一个 \emph{双模拟（bisimulation）}。当 $M_1$ 与 $M_2$ 之间存在一个连接世界 $w_1$ 与 $w_2$ 的双模拟时，我们也可以写作 $\tuple{M_1, w_1} \leftrightarroweq \tuple{M_2, w_2}$，并把 $\tuple{M_1, w_1}$ 与 $\tuple{M_2, w_2}$ 称为 \emph{双模拟等价的（bisimilar）}。
\end{defn}

双模拟关系中的不同条款保证了不同的事情。条款 1 保证双模拟等价的世界满足相同的无模态公式，因为它保证了对所有!!{propositional variable}的一致。另外两个条款——有时分别被称为「forth」和「back」——保证可及关系具有相同的结构。

\begin{thm}
如果 $\tuple{M_1, w_1} \leftrightarroweq \tuple{M_2, w_2}$，则对每个!!{formula}~$!A$，我们有 $\mSat{M_1}{!A}[w_1]$ 当且仅当 $\mSat{M_2}{!A}[w_2]$。
\end{thm}

\begin{figure}
  \begin{center}
    \begin{tikzpicture}[modal]
      \node[world] (w1) {$w_1$};
      \node[world] (w2) [above left=of w1]{$w_2$};
      \node[world] (w3) [above right=of w1] {$w_3$};
      \draw[<->] (w1) to node {$a$} (w2);
      \draw[->,loop below] (w1) to node {$a$} (w1);
      \draw[<->] (w1) to [swap] node {$a$} (w3);
      \draw[->,loop above] (w2) to node {$a$} (w2) ;
      \draw[->,loop above] (w3) to node {$a$} (w3) ;

      \node[world](v1)[right of=w1, xshift=1.5in]{$v_1$};
      \node[world](v2)[above of=v1]{$v_2$};
      \draw[<->] (v1) to node {$a$} (v2);
      \draw[->,loop below] (v1) to node {$a$} (v1);
      \draw[->,loop above] (v2) to node {$a$} (v2) ;

      \draw[-,dotted] (w1) to (v1) ;
      \draw[-,dotted,bend left=45] (w2) to (v2) ;
      \draw[-,dotted,bend left=30] (w3) to (v2) ;
    \end{tikzpicture}
  \end{center}
  \caption{两个双模拟等价的模型。}
  \ollabel{fig:bisimilar}
\end{figure}

尽管 \olref{fig:bisimilar} 中画出的两个模型彼此并不完全相同，但存在一个连接世界 $w_1$ 与~$v_1$ 的双模拟。这个双模拟也会把 $w_2$、$w_3$ 二者同时与~$v_2$ 连接起来，其思想在于：我们的模态语言中没有任何可表达的东西能真正区分它们。不过，如果 $w_2$ 与~$w_3$ 满足不同的!!{propositional variable}，情况就会不同。

\end{document}
